\part{电磁学}
\chapter{静电场}
\section{电荷~库仑定律}

\subsection{电荷}

1.\textbf{电荷量}：表示物体所带电荷多寡的物理量。

2.物体所带电荷只有正电荷和负电荷,同号电荷相互排斥，异号电荷相互吸引。

\subsection{电荷守恒定律}

1.\textbf{起电}：通过某种作用，使该物体内电子不足或过多而呈带电状态。

2.\textbf{电荷守恒定律}：在一个与外界没有电荷交换的系统内，无论经过怎样的物理过程，系统内正、负电荷的代数和总保持不变。

3.电荷是相对论不变量，即电荷量与运动无关。

\subsection{电荷的量子化}

电荷量这种只能取分立的、不连续量值的性质，称为电荷量子化。

一般近似取$e=1.602176634×10^{-19}$

\subsection{库仑定律}

1.库仑扭秤实验总结出了\textbf{库仑定律}：两个静止点电荷之间相互作用力的大小与这两个点电荷的电荷$q_{1}$和$q_{2}$的乘积成正比，而与这两个点电荷之间的距离$r_{12}$的二次方成反比，作用力的方向，沿着这两个点电荷的连线。

$$\mathbf{F}=k\frac{q_{1}q_{2}}{{r_{12}}^{2}}$$
$$k=\frac{1}{4\pi\varepsilon_{0}}$$
则$$\mathbf{F_{12}}=\mathbf{F_{21}}=\frac{1}{4\pi\varepsilon_{0}}\frac{q_{1}q_{2}}{{r_{12}}^{2}}\mathbf{e_{r_{12}}} $$

2.适用条件：真空、静止、点电荷

\subsection{静电力的叠加原理}

1.\textbf{静电力叠加原理}：当空间有两个以上的点电荷时，作用在某一点电荷上的总静电力等于其他各点电荷单独存在时，对该点电荷所施静电力的矢量和。

$$\mathbf{F}=\sum_{i=1}^{n}\mathbf{F_{i}}=\frac{1}{4\pi\varepsilon_{0}}\sum_{i=1}^{n}\frac{qq_{i}}{{r_{i}}^{2}}\mathbf{e_{r_{i}}}$$

2.库伦定律只适用于点电荷,将带电体看作是由许多电荷元组成的，根据库仑定律和叠加原理，则可求解带电体之间的相互作用力。
